\section{Legality and recursive independence}
\label{sec:legality}

The auxiliary inequalities are weak, but the public answers are strict.
Nonempty auxiliary counts alone would not establish legality: a strict
directed cycle can still admit a constant weak assignment. We instead prove
that the public order remains acyclic and that the recursive restrictions
can always be realized together.

\begin{lemma}[Effective support]
\label{lem:support}
The interval system is feasible exactly when $\ell_x\leq r_x$ for every item.
In a feasible system every integer in $[\ell_x,r_x]$ is attained by $X_x$.
Along an active edge $x\to y$, $\ell_x\leq\ell_y$ and $r_x\leq r_y$.
\end{lemma}
\begin{proof}
Every assignment respects propagated bounds. Conversely the lower-endpoint
vector $(\ell_x)_x$ respects every interval and inequality if
$\ell_x\leq r_x$. To obtain $X_x=t$, add the bound $[t,t]$ at $x$ and
propagate again. No predecessor lower bound exceeds $t$, and no successor
upper bound is below $t$, so the result is feasible. The final monotonicity
statements follow from the definitions of the propagated endpoints.
\end{proof}

\begin{lemma}[Interior fixing and endpoint insertion]
\label{lem:split}
Fixing $X_x=c$ with $\ell_x<c<r_x$ leaves at least two feasible positions for
every other free item. The item $x$ is publicly incomparable with all old
items fixed at $c$.
If a selected position is an endpoint, the insertion of
Section~\ref{sec:move} makes its new copy interior and retains at least as
many assignments as the original endpoint fiber.
\end{lemma}
\begin{proof}
Only predecessors of $x$ can acquire a smaller upper bound after interior
fixing; their lower bounds are at most $\ell_x<c$. Only successors can
acquire a larger lower bound, and their upper bounds are at least $r_x>c$.
Thus no other free interval becomes a singleton.
An old fixed item $p$ at $c$ with $p<_P x$ would imply $X_x\geq c$
in every assignment, contradicting attainment of $\ell_x<c$.
The case $x<_P p$ is symmetric.

For a lower endpoint $c$, take an old assignment with $X_x=c$.
Raise $x$ and exactly its successors currently equal to $c$ to the new
position $c^+$. No such successor is fixed, since a fixed successor at
$c$ would give $r_x\leq c$. All inequalities survive, and every changed
artificial interval admits the new position. Collapsing $c^+$ back to $c$
recovers the original assignment, proving injectivity. Assignments attaining
the old lower and upper endpoints still exist before fixing; hence $c^+$
is interior. Reflect this argument for an upper endpoint.

It is essential that comparison-induced bounds are recomputed rather than
enlarged as artificial bounds. A comparison $x<p$ with $p$ fixed at the old
identity $c$ still implies $X_x\leq c$, not $X_x\leq c^+$.
\end{proof}

\begin{lemma}[Independent admissions]
\label{lem:birth}
The $j$th admission to a child multiplies its count by at least $j$.
It changes neither the effective intervals nor the probability distribution
of its old free items.
\end{lemma}
\begin{proof}
Lemma~\ref{lem:split} makes the incoming item publicly incomparable with
every old item in that child. It therefore introduces no inequality there.
Every old assignment has exactly as many extensions as there are positions
in the incoming item's full interval, at least $j$. Appending a position
does not enlarge any old interval, so this extension count is constant
over old assignments. The assertion remains true if the child has answered
comparisons between successive admissions.
\end{proof}

\begin{lemma}[Safe routing]
\label{lem:routing}
For a publicly incomparable pair $a,b$ not both fixed at the same position,
either both answers are safe, or a safe answer retains every old assignment.
Delegation through a common fixed position changes only the count at the
eventual stopping node and introduces no extra comparison into another child.
Every resulting recursive state admits a strict total order extending $P$.
\end{lemma}
\begin{proof}
Suppose adding $a<b$ creates a new public relation $i<j$ between two old
items fixed at position $c$. There were old public paths
$i\leq_P a$ and $b\leq_P j$. Consequently every old assignment satisfies
\begin{equation}
                         X_b\leq c\leq X_a.
\label{eq:reverse}
\end{equation}
The reverse weak answer retains everything. If both directions have such
an obstruction, the two inequalities force $X_a=X_b=c$ throughout the
old family. By Lemma~\ref{lem:support} and the two-position invariant,
both items were already fixed at that position: the delegation case.

Suppose instead that $a<b$ newly confines a free item $w$ to $c$.
Both bounds of $w$ cannot have changed, since that would require an old
path $b\leq_P w\leq_P a$, contrary to incomparability.
If the lower bound rose, $b\leq_P w$, the old upper bound of $w$ was
$c$, and the old lower bound of $a$ was $c$. This again gives
\eqref{eq:reverse}. There was an old assignment with $X_w<c$, so there
was a strict witness $X_b<X_a$. That witness excludes an inner-fixed-item
obstruction to the reverse answer, since such an obstruction would force
$X_a\leq X_b$ in every old assignment. The reverse cannot create a
singleton either, because it retains the whole old family. A falling upper
bound is symmetric. If a weak answer is infeasible, its reverse is strict
in every old assignment, giving the same conclusion.

If both items are fixed at $c$, either weak answer retains every parent
assignment. A newly acquired relation between items in this child is
precisely a consequence of its requested comparison: its old path pieces
are already in the child's public restriction. If another child at $d$
acquired a relation, those pieces would imply $c\leq d\leq c$, so
$d=c$. Thus delegation affects no other child.

Finally choose an admissible assignment at each node. Order different
positions increasingly. Within a tied position, extend the restriction of
$P$, respecting the order recursively chosen for its fixed child items.
Any path through free items imposes only a relation already in that public
restriction. A total order extending a restriction of a finite partial order
can be extended to the larger set, so these choices can be combined.
Publicly incomparable queries never create a directed cycle; the safety
tests ensure consistency with the existing child restrictions. Recursion
therefore yields a strict total order on the original distinct items.
\end{proof}

The public order is the sorter's only information. Private interval cuts and
fixings do not become extra answers. Consequently a correct sorter can stop
only when $P$ is total, even if it knows the adversary's program: it must also
be correct on every other input giving the same transcript.
